\documentclass[11pt,a4paper]{article}

% ----------------------------------------------------------------------------
% Packages
% ----------------------------------------------------------------------------
\usepackage[utf8]{inputenc}
\usepackage[T1]{fontenc}
\usepackage{lmodern}
\usepackage{microtype}
\usepackage[margin=2.4cm]{geometry}
\usepackage{amsmath,amssymb}
\usepackage{booktabs}
\usepackage{listings}
\usepackage{xcolor}
\usepackage{tcolorbox}
\tcbuselibrary{skins,breakable}
\usepackage{enumitem}
\usepackage{graphicx}
\usepackage[hidelinks]{hyperref}
\usepackage{titlesec}

% ----------------------------------------------------------------------------
% Colours & code style
% ----------------------------------------------------------------------------
\definecolor{codebg}{HTML}{F6F8FA}
\definecolor{codekw}{HTML}{0033B3}
\definecolor{codestr}{HTML}{067D17}
\definecolor{codecmt}{HTML}{8C8C8C}
\definecolor{rule}{HTML}{1F4E79}
\definecolor{accent}{HTML}{8A2BE2}

\lstdefinestyle{py}{
  language=Python,
  backgroundcolor=\color{codebg},
  basicstyle=\ttfamily\small,
  keywordstyle=\color{codekw}\bfseries,
  stringstyle=\color{codestr},
  commentstyle=\color{codecmt}\itshape,
  numberstyle=\tiny\color{codecmt},
  numbers=left,
  numbersep=8pt,
  showstringspaces=false,
  breaklines=true,
  frame=single,
  framerule=0pt,
  rulecolor=\color{codebg},
  xleftmargin=12pt,
  aboveskip=8pt,belowskip=8pt,
  literate={->}{{$\rightarrow$}}1 {>=}{{$\geq$}}1 {<=}{{$\leq$}}1
}
\lstset{style=py}

% ----------------------------------------------------------------------------
% Teaching boxes
% ----------------------------------------------------------------------------
\newtcolorbox{keyidea}[1][]{colback=accent!6,colframe=accent!70,
  fonttitle=\bfseries,title=Key idea,breakable,#1}
\newtcolorbox{pitfall}[1][]{colback=red!4,colframe=red!55,
  fonttitle=\bfseries,title=Common pitfall,breakable,#1}
\newtcolorbox{specbox}[1][]{colback=rule!5,colframe=rule!70,
  fonttitle=\bfseries,title=What the specification says,breakable,#1}
\newtcolorbox{defbox}[2][]{colback=black!3,colframe=black!55,
  fonttitle=\bfseries,title=Definition: #2,breakable,#1}

\titleformat{\section}{\Large\bfseries\color{rule}}{\thesection}{0.6em}{}
\titleformat{\subsection}{\large\bfseries\color{rule!85}}{\thesubsection}{0.6em}{}

\newcommand{\code}[1]{\texttt{#1}}

% ----------------------------------------------------------------------------
\title{\vspace{-1.2cm}\textbf{Engineering for Conformance}\\[2pt]
  \large A Case Study in Bringing an Emulator into Line with a Real-World API\\[6pt]
  \normalsize\textit{The Philips Hue API and the diyHue Bridge Emulator}}
\author{Software Engineering Teaching Notes}
\date{\today}

\begin{document}
\maketitle

\begin{abstract}
\noindent
This chapter is a worked example of a recurring task in professional software
engineering: making one program faithfully imitate another whose behaviour is
defined by a \emph{specification} rather than by shared source code. Our subject
is \textsf{diyHue}, an open-source emulator of the Philips Hue smart-lighting
bridge. We start from a real, automated measurement showing the emulator is only
$\approx 33\%$ conformant with the Hue API, then repair nine subsystems and lift
conformance to $100\%$ of an automated check suite (and a full test suite from
$44$ passing / $38$ failing to $110$ passing / $0$ failing). Each repair is
presented as a small, self-contained lesson: the rule being violated, the bug,
the fix, and the general principle it illustrates. The recurring themes are
\emph{input validation}, \emph{representation invariants}, \emph{total
functions}, \emph{idempotent reconciliation}, and \emph{executable
specifications}.
\end{abstract}

\begin{keyidea}
A specification is a \emph{contract}. When you cannot see the other party's
code, the contract is all you have. ``It works on my device'' is not the same
as ``it obeys the contract'', and only the second one makes independent clients
interoperate.
\end{keyidea}

\section*{Learning objectives}
\addcontentsline{toc}{section}{Learning objectives}
After studying this chapter you should be able to:
\begin{itemize}[noitemsep]
  \item explain the difference between a \emph{test} and a \emph{conformance
        scorecard}, and why both are useful;
  \item identify the classic defect classes---missing input validation,
        broken representation invariants, partial (non-total) functions,
        unsafe exception handling, and identity-vs-name confusion;
  \item apply modular arithmetic and clamping to validate bounded inputs;
  \item reason about ownership graphs and idempotent reconciliation;
  \item recognise why catching \code{Exception} can hide control-flow bugs in
        Python.
\end{itemize}

\tableofcontents
\bigskip

% ============================================================================
\section{Background: two APIs and one emulator}
% ============================================================================

A Philips Hue \emph{bridge} is a small Linux device that speaks ZigBee to the
lamps and HTTP to phones and home-automation software. Two generations of HTTP
interface coexist:

\begin{itemize}
  \item the \textbf{v1 (legacy) REST API} under \code{/api/<user>/\ldots},
        whose light state uses integer fields such as \code{bri}
        ($1$--$254$), \code{ct} ($153$--$500$ mireds), and \code{hue}
        ($0$--$65535$, wrapping); and
  \item the \textbf{CLIP v2 API} under \code{/clip/v2/resource/\ldots}, a
        resource-oriented model where each physical product is a
        \code{device} resource that \emph{owns} a set of \emph{service}
        resources (a \code{light}, a \code{button}, a \code{zigbee\_con\-nec\-tiv\-ity},
        \ldots). Brightness here is a floating-point \emph{percentage}
        ($0$--$100$).
\end{itemize}

\textsf{diyHue} re-implements both, plus a Server-Sent-Events (SSE) stream and a
DTLS ``entertainment'' channel, on top of arbitrary hardware (ZigBee, MQTT,
WLED, ESPHome, \ldots). Because real Hue apps were written and tested against a
\emph{real} bridge, the emulator only works if it reproduces the real bridge's
observable contract precisely.

\begin{defbox}{conformance}
\emph{Conformance} is the degree to which an implementation's observable
behaviour matches a specification. It is observable: it says nothing about
\emph{how} the program is built, only about what a client can see.
\end{defbox}

% ============================================================================
\section{Method: measure before you mend}
% ============================================================================

Before changing code we built an automated \textbf{conformance scorecard}: a
program that pokes the emulator and checks each spec rule, then prints a single
weighted percentage. Crucially, a failing check does \emph{not} abort the run;
it is recorded and reported, so the scorecard always produces a number you can
watch rise as you work.

\begin{keyidea}
\textbf{Make the invisible measurable.} ``Improve compliance'' is not
actionable. ``Raise the weighted conformance score from $33\%$ to $95\%$'' is.
A good metric turns a vague goal into a gradient you can descend.
\end{keyidea}

Each check carries a severity weight ($\textsc{high}=3$, $\textsc{med}=2$,
$\textsc{low}=1$), and the overall score is
\[
  \text{score} \;=\;
  100 \times
  \frac{\sum_{c\,\in\,\text{passing}} w_c}{\sum_{c\,\in\,\text{all}} w_c}.
\]
Weighting by severity means fixing one crash counts for more than polishing
three cosmetic fields---the metric encodes our priorities.

\begin{table}[h]
\centering
\caption{Conformance before and after the work in this chapter.}
\begin{tabular}{lcc}
\toprule
 & \textbf{Before} & \textbf{After}\\
\midrule
Weighted conformance score & $32.8\%$ (grade E) & $\mathbf{100\%}$ (grade A)\\
Scorecard checks passing    & $9/27$            & $\mathbf{27/27}$\\
Full unit/contract suite    & $44$ pass / $38$ fail & $\mathbf{110}$ pass / $\mathbf{0}$ fail\\
\bottomrule
\end{tabular}
\end{table}

The scorecard is itself written with \textsf{pytest}. A non-conformant probe is
marked \code{xfail} (``expected failure''), so the suite stays green while still
recording the gap; as a bug is fixed the probe silently flips to a pass and the
score rises with no edit to the check. A final \emph{regression gate} asserts
the score never falls below a floor, so a future change that breaks conformance
fails the build.

\begin{defbox}{executable specification}
An \emph{executable specification} is a checkable encoding of the rules a
program must obey. Unlike prose, it cannot drift out of date silently: if the
rule and the code disagree, the test run tells you.
\end{defbox}

% ============================================================================
\section{Case studies}
% ============================================================================

We now walk through the nine repairs. Each is a miniature lesson; read the
``Key idea'' boxes even if you skim the code.

% ----------------------------------------------------------------------------
\subsection{Input validation: clamping and wrapping}
% ----------------------------------------------------------------------------

\begin{specbox}
A v1 light command must respect field ranges: \code{bri}$\in[1,254]$,
\code{ct}$\in[153,500]$, \code{sat}$\in[0,254]$, and \code{hue} is a $16$-bit
value that \emph{wraps} modulo $65536$. \code{bri}=$0$ is \emph{not} valid
(``off'' is \code{on}=\textsf{false}, not brightness zero)~\cite{huedev,openhue}.
\end{specbox}

The emulator wrote whatever the client sent straight into its state, so
\code{\{"bri":999\}} was stored verbatim. We added a single validation function
applied on every direct state write:

\begin{lstlisting}
V1_STATE_RANGES = {"bri": (1, 254), "ct": (153, 500), "sat": (0, 254)}

def clampV1State(state):
    """Clamp/wrap v1 light-state values into spec-valid ranges, in place."""
    for key, (lo, hi) in V1_STATE_RANGES.items():
        if key in state and isinstance(state[key], (int, float)) \
                and not isinstance(state[key], bool):
            state[key] = max(lo, min(hi, int(state[key])))
    if "hue" in state and isinstance(state["hue"], (int, float)):
        state["hue"] = int(state["hue"]) % 65536      # wrap, do not clamp
    ...
    return state
\end{lstlisting}

Two details deserve attention. First, \code{hue} is \emph{wrapped}, not clamped:
the original code did \code{state["hue"] -= 65535}, which is both off-by-one
(the period is $65536$, not $65535$) and wrong for large increments. The
correct operation is a single modulo:
\[
  \text{hue}' \;=\; (\text{hue} + \Delta) \bmod 65536 .
\]
Second, \code{bool} is a subclass of \code{int} in Python, so
\code{isinstance(True, int)} is \code{True}~\cite{pystdtypes}; we exclude booleans explicitly so
that \code{on=True} is never mistaken for a numeric field.

\begin{pitfall}
\textbf{Trusting the client.} Never assume input is in range just because a
well-behaved client would send valid values. A different client---or a
malicious one---will not. Validate at the boundary, once, in a single place you
can point to.
\end{pitfall}

% ----------------------------------------------------------------------------
\subsection{Representation invariants: the response envelope}
% ----------------------------------------------------------------------------

\begin{specbox}
Every CLIP v2 response is an object with \emph{both} an \code{errors} array and
a \code{data} array---even on error (then \code{data} is empty) and even on
success (then \code{errors} is empty)~\cite{openhue,huev2}.
\end{specbox}

The emulator's \code{PUT} handler returned only \code{data}; its error path
returned only \code{errors} and even \code{del}eted the \code{data} key. A
client that unconditionally reads \code{response["data"]} would crash. The fix
is to make the shape an \emph{invariant}---always present, regardless of
outcome:

\begin{lstlisting}
# success
response = {"errors": [], "data": [{"rid": resourceid, "rtype": resource}]}
# error
return {"errors": [{"description": f"...{resource}"}], "data": []}, 500
\end{lstlisting}

\begin{defbox}{representation invariant}
A \emph{representation invariant} is a property that every value of a type must
satisfy. ``A v2 response always has \code{errors} and \code{data} arrays'' is
one. Enforcing invariants at the point of construction means clients never have
to defend against malformed values.
\end{defbox}

A sibling bug lived in the ``get one button by id'' path, which returned
\code{data:[object.getButtons()]}---a \emph{list inside a list}, because
\code{getButtons()} already returns a list. The repair flattens the result and,
because a single button is a sub-resource of a sensor device, resolves it by
searching for the one matching id:

\begin{lstlisting}
if resource in ["button", "relative_rotary"]:
    method = "getButtons" if resource == "button" else "getRotary"
    items = [item for sensor in bridgeConfig["sensors"].values()
             for item in getattr(sensor, method)()
             if item["id"] == resourceid]
    return {"errors": [], "data": items}
\end{lstlisting}

% ----------------------------------------------------------------------------
\subsection{Data modelling: the device/service ownership graph}
% ----------------------------------------------------------------------------

The v2 model is a small graph: a \code{room} has \emph{children} (the devices in
it) and \emph{services} (its \code{grouped\_light}); a \code{device} owns its
services; a \code{grouped\_light} is owned by its room or zone~\cite{openhue,huev2}. Three bugs lived
here.

\paragraph{(a) Missing \texttt{owner}.} Every service and device must carry an
\code{owner} reference \code{\{rid, rtype\}}. Devices omitted it, failing schema
validation. We added it at each device's single construction point.

\paragraph{(b) Non-light children were dropped.} A room could contain switches
and sensors, but \code{getV2Room} only emitted its lights. We gave \code{Group}
a \code{device\_children} list that is initialised from stored data (as a
\emph{copy}, so later mutation of the source cannot leak in), rendered into the
children array, de-duplicated against the lights, and persisted on save:

\begin{lstlisting}
self.device_children = list(data["device_children"]) \
    if "device_children" in data else []
...
seen = set()
for light in self.lights:
    if light():
        rid = light().getDevice()["id"]
        if rid not in seen:
            seen.add(rid); result["children"].append({"rid": rid, "rtype": "device"})
for rid in self.device_children:            # non-light devices, deduped
    if rid not in seen:
        seen.add(rid); result["children"].append({"rid": rid, "rtype": "device"})
\end{lstlisting}

\paragraph{(c) Wrong owner on \texttt{grouped\_light}.} It claimed to be owned by
a \emph{device}; the spec says a grouped light is owned by its room/zone (or, for
group $0$, by \code{bridge\_home}). We compute the owner from the group's type.

\begin{keyidea}
\textbf{Model identities, not names.} A reference points at an \code{id}, never
at a display name. Names change; identities do not. Half the bugs in this
section are a node pointing at the wrong identity.
\end{keyidea}

% ----------------------------------------------------------------------------
\subsection{Event semantics: batching, cascades, and an import trap}
% ----------------------------------------------------------------------------

\begin{specbox}
A single light change is delivered as \emph{one} SSE event whose \code{data}
array contains every affected resource: the light, its device, and the
\code{grouped\_light} rollups for every owning room/zone and for
\code{bridge\_home} (the ``all lights'' group)~\cite{whatwg-sse,huev2}.
\end{specbox}

The emulator emitted a \emph{separate} event per resource and never sent the
group rollups, so a dashboard tracking a room's aggregate brightness silently
went stale. We rewrote the emitter to build one batched event and append the
rollups for every group containing the light.

This subsection carries a subtle and very instructive Python lesson. To find the
owning groups, the light needs the global configuration. A naive
\code{import configManager} \emph{inside} the function re-runs that module's
top-level code, which parses command-line arguments and calls
\code{sys.exit()} when run under the test harness. That raises
\code{SystemExit}---which is \textbf{not} a subclass of \code{Exception}~\cite{pyexc}:

\[
  \texttt{BaseException} \supset \{\texttt{SystemExit},\ \texttt{KeyboardInterrupt},\ \texttt{Exception} \supset \texttt{KeyError}, \ldots\}
\]

so a surrounding \code{except Exception} does \emph{not} catch it and the whole
test crashes. The fix never triggers a fresh import; it reads the
\emph{already-imported} module out of the module cache:

\begin{lstlisting}
import sys
groups = {}
cfgMod = sys.modules.get("configManager")   # do NOT `import` here
if cfgMod is not None:
    try:
        groups = cfgMod.bridgeConfig.yaml_config.get("groups", {})
    except Exception:
        groups = {}
\end{lstlisting}

\begin{pitfall}
\textbf{\code{import} runs code.} In Python, importing a module executes its
top level. Importing for a side effect deep inside a hot function can re-run
that code in a surprising context. And remember: \code{except Exception} does
\emph{not} catch \code{SystemExit} or \code{KeyboardInterrupt}---those are
\code{BaseException}s by design~\cite{pyexc}, so a stray \code{sys.exit()} sails right past
your handler.
\end{pitfall}

% ----------------------------------------------------------------------------
\subsection{Total functions: the rules-engine fall-through}
% ----------------------------------------------------------------------------

The v1 rules engine evaluates conditions of the form
\code{\{address, operator, value\}}. The dispatch was a chain of
\code{if/elif} on the operator, with \emph{no final \code{else}}. An operator
the engine did not recognise---\code{"not in"}, \code{"stable"}---matched no
branch, so the loop simply fell through and the condition was treated as
\emph{satisfied}. A rule that should never fire could therefore fire.

\begin{keyidea}
\textbf{Make dispatch total.} A function that switches on a value should handle
\emph{every} value, including the unexpected ones. The safe default for an
unknown operator is ``condition not met'', never ``silently true''. In practice
this means: always write the final \code{else}.
\end{keyidea}

\begin{lstlisting}
elif condition["operator"] == "not in":
    ...                              # genuinely evaluate the time window
elif condition["operator"] in ("stable", "not stable"):
    return [False, 0]                # cannot evaluate -> do not fire
else:
    logging.debug("unsupported operator " + str(condition["operator"]))
    return [False, 0]                # unknown -> safe default
\end{lstlisting}

% ----------------------------------------------------------------------------
\subsection{Robustness: a crash hiding behind a missing key}
% ----------------------------------------------------------------------------

Creating a v2 \code{behavior\_instance} with a perfectly valid payload returned
HTTP $500$. The cause was one line that assumed an optional field always
exists:
\begin{lstlisting}
for resource in self.configuration["where"]:        # KeyError if absent
\end{lstlisting}
The repair is a one-character class of fix---ask for the key \emph{defensively}:
\begin{lstlisting}
for resource in self.configuration.get("where", []):
\end{lstlisting}

\begin{pitfall}
\code{d[k]} demands that \code{k} exists; \code{d.get(k, default)} tolerates its
absence. Use indexing only when absence is genuinely a bug worth crashing for.
For optional fields, supply a default.
\end{pitfall}

% ----------------------------------------------------------------------------
\subsection{Resource management: rate limiting}
% ----------------------------------------------------------------------------

\begin{specbox}
A real bridge processes only about \emph{ten} light commands per second; beyond
that it returns HTTP $429$ (``Too Many Requests'')~\cite{rfc6585,huedev}. Apps are written to expect
this back-pressure.
\end{specbox}

The emulator applied every command immediately, so a client tuned to the real
bridge could overwhelm slow hardware. We added a \emph{sliding-window} limiter:
keep the timestamps of recent commands, drop those older than one second, and
reject once the window is full.

\begin{lstlisting}
_RATE_LIMITS = {"light": 10}
_rate_state = {}

def rateLimited(resource):
    limit = _RATE_LIMITS.get(resource)
    if not limit:
        return False
    now = monotonic()
    window = _rate_state.setdefault(resource, [])
    window[:] = [t for t in window if t > now - 1.0]   # forget old events
    if len(window) >= limit:
        return True
    window.append(now)
    return False
\end{lstlisting}

\begin{defbox}{sliding-window rate limit}
A \emph{sliding-window} limiter permits at most $N$ events in any trailing time
window of length $T$. It is implemented by remembering recent event times and
discarding those older than $T$ before each decision. Note the use of a
\emph{monotonic} clock, which never jumps backwards when the system time is~\cite{pytime}
adjusted.
\end{defbox}

% ----------------------------------------------------------------------------
\subsection{Idempotent reconciliation: de-duplicating MQTT sensors}
% ----------------------------------------------------------------------------

When a ZigBee device is renamed in the Zigbee2MQTT layer, the bridge re-announces
it under a new friendly name. The emulator keyed sensors by \emph{name}, so a
rename created a \emph{second} copy instead of updating the first---and the two
copies fought over the device. The real identity of a ZigBee device is not its
name but its \code{ieeeAddr} (its hardware address). The fix groups sensors by
that stable key, keeps exactly one per logical type, and renames in place:

\begin{lstlisting}
def getSensorsByIeeeAddr(ieeeAddr):
    return [s for s in bridgeConfig["sensors"].values()
            if getattr(s, "protocol", None) == "mqtt"
            and s.protocol_cfg.get("ieeeAddr") == ieeeAddr]
\end{lstlisting}

\begin{keyidea}
\textbf{Reconcile, don't duplicate.} Processing the same announcement twice
should leave the system in the same state as processing it once. This property
is called \emph{idempotence}, and the way to get it is to key your records by a
stable identity and \emph{update} rather than \emph{insert}.
\end{keyidea}

% ----------------------------------------------------------------------------
\subsection{Schema-driven dispatch: the tap-dial automation}
% ----------------------------------------------------------------------------

The most involved repair re-implemented the automation engine for the RDM002
``tap dial'' (four buttons plus a rotary). The button-event integer encodes both
\emph{which} button and \emph{what} happened:
\[
  \underbrace{1}_{\text{button }1}\,\underbrace{002}_{\text{short release}}
\]
i.e. \code{button = event // 1000} and \code{action = event \% 1000}. We map the
trailing digit to a named action through a small table---clearer and more
total than a chain of \code{if}s:

\begin{lstlisting}
BUTTON_ACTIONS = {0: "on_initial_press", 1: "on_repeat",
                  2: "on_short_release", 3: "on_long_press"}
...
buttonKey = "button" + str(buttonevent // 1000)
actionKey = BUTTON_ACTIONS.get(buttonevent % 1000)
\end{lstlisting}

The rotary dial is a separate logical sensor whose automation is configured on
its \emph{parent} switch; matching therefore checks both the device's own id and
its \code{parent\_id\_v2}. The new schema lives alongside the legacy one (we
fall back to the old format when the new keys are absent), so existing stored
automations keep working---a small but important act of
\emph{backward compatibility}.

% ============================================================================
\section{Verification: the scorecard as a safety net}
% ============================================================================

Every repair was confirmed by re-running the scorecard and watching the number
climb: $32.8\% \to 62.3\% \to 96.7\% \to 100\%$. Because the checks are
executable, they double as regression protection: the gate
\code{test\_scorecard\_meets\_baseline} now fails the build if conformance drops
below $95\%$.

\begin{keyidea}
\textbf{A fix without a test is a rumour.} Each defect in this chapter is now
pinned by a check that would have caught it. The scorecard is not scaffolding to
be discarded---it is the executable definition of ``done''.
\end{keyidea}

% ============================================================================
\section{Recurring principles}
% ============================================================================
\begin{enumerate}[leftmargin=2em]
  \item \textbf{Measure first.} Turn ``be more compliant'' into a number, then
        make the number go up.
  \item \textbf{Validate at the boundary.} Clamp and wrap untrusted input once,
        in one place.
  \item \textbf{Enforce invariants at construction.} If every response has the
        same shape, no client has to defend against malformed ones.
  \item \textbf{Make functions total.} Handle every case; the safe default for
        the unknown case is rarely ``succeed''.
  \item \textbf{Key by identity, not by name.} Stable identifiers make
        reconciliation idempotent.
  \item \textbf{Know your language.} \code{import} runs code; \code{SystemExit}
        is not an \code{Exception}; \code{bool} is an \code{int}.
\end{enumerate}

% ============================================================================
\section*{Exercises}
% ============================================================================
\addcontentsline{toc}{section}{Exercises}
\begin{enumerate}[leftmargin=2em]
  \item \textbf{(Wrapping.)} Show that \code{(hue + }$\Delta$\code{) \% 65536}
        gives the correct result for $\Delta = -90000$, whereas a single
        \code{+= 65536} correction does not. For which range of $\Delta$ does
        the single-correction approach work?
  \item \textbf{(Invariants.)} A teammate proposes returning bare \code{[]} for
        an empty collection instead of \code{\{"errors":[],"data":[]\}}.
        Give a concrete client that breaks, and state the invariant being
        violated.
  \item \textbf{(Totality.)} The rules engine fell through on unknown operators.
        Write a property-based test that feeds random operator strings and
        asserts the condition is never reported as satisfied unless the operator
        is one of the supported set.
  \item \textbf{(Rate limiting.)} The limiter uses \code{monotonic()} rather
        than \code{time()}. Construct a scenario (an NTP step, a daylight-saving
        change) where using wall-clock time would let a burst through or block a
        legitimate command.
  \item \textbf{(Idempotence.)} Prove that processing the same
        \code{bridge/devices} announcement twice leaves the sensor table
        identical to processing it once, given the de-duplication rule
        ``one sensor per type, keyed by \code{ieeeAddr}''.
  \item \textbf{(Design.)} The tap-dial decodes \code{event // 1000} and
        \code{event \% 1000}. What goes wrong if a future device emits button
        index $10$? Propose an encoding that does not have this ceiling and
        discuss its compatibility cost.
\end{enumerate}

% ============================================================================
\begin{thebibliography}{9}
\addcontentsline{toc}{section}{References}

\bibitem{huedev} Signify / Philips Hue. \emph{Hue Developer Documentation --- CLIP
API v2 and legacy API, core concepts (resource model, rate limits, light state
ranges)}. \url{https://developers.meethue.com/develop/hue-api-v2/}. (Access
requires a free developer account; the machine-readable transcription in
\cite{openhue} mirrors its field names and enums.)

\bibitem{openhue} OpenHue. \emph{openhue-api --- an OpenAPI 3 specification of the
Philips Hue CLIP v2 API}. \url{https://github.com/openhue/openhue-api}. Source of
the device/service composition model, the \texttt{\{errors,data\}} envelope, and
the \texttt{dimming.brightness} $0$--$100$ percentage.

\bibitem{rfc6585} R. Fielding and M. Nottingham. \emph{RFC 6585: Additional HTTP
Status Codes}. IETF, April 2012. \url{https://datatracker.ietf.org/doc/html/rfc6585}.
Defines \texttt{429 Too Many Requests} and the \texttt{Retry-After} header.

\bibitem{whatwg-sse} WHATWG. \emph{HTML Living Standard, \S9.2 Server-sent events}.
\url{https://html.spec.whatwg.org/multipage/server-sent-events.html}. The
\texttt{text/event-stream} wire format (\texttt{id:}/\texttt{data:} fields) and
the \texttt{EventSource} API.

\bibitem{pyexc} Python Software Foundation. \emph{Python Language Reference ---
Built-in Exceptions: the exception hierarchy}.
\url{https://docs.python.org/3/library/exceptions.html}. \texttt{SystemExit} and
\texttt{KeyboardInterrupt} derive from \texttt{BaseException}, not
\texttt{Exception}, so \texttt{except Exception} does not catch them.

\bibitem{pystdtypes} Python Software Foundation. \emph{The Python Standard Library
--- Numeric Types and Boolean Values}.
\url{https://docs.python.org/3/library/stdtypes.html\#boolean-type-bool}.
\texttt{bool} is a subclass of \texttt{int} (\texttt{True}=1, \texttt{False}=0;
see also PEP 285).

\bibitem{pytime} Python Software Foundation. \emph{The Python Standard Library ---
\texttt{time.monotonic()}}. \url{https://docs.python.org/3/library/time.html\#time.monotonic}.
A clock that cannot go backwards and is unaffected by system-clock adjustments
(rationale in PEP 418).

\bibitem{diyhue} The diyHue project. \emph{diyHue --- a Philips Hue bridge
emulator (source repository)}. \url{https://github.com/diyhue/diyHue}. Primary
source for the emulator behaviour discussed here; specific defects cite
\texttt{BridgeEmulator/\ldots} by file and line.

\bibitem{companion} Companion artifacts in this repository:
\texttt{docs/spec-compliance-audit.md} (the full defect inventory),
\texttt{docs/real-bridge-research.md} (measured behaviour of the genuine BSB002
bridge), and \texttt{tests/test\_spec\_compliance.py} (the executable scorecard).

\end{thebibliography}

\vfill
\noindent\rule{\linewidth}{0.4pt}\\
\footnotesize These notes accompany the change set on the
\texttt{spec-compliance-fixes} branch of \textsf{diyHue}. The companion
documents \texttt{spec-compliance-audit.md} (the full defect list) and
\texttt{real-bridge-research.md} (how the real bridge behaves) provide the
underlying evidence.

\end{document}
